% resumo em português
\setlength{\absparsep}{18pt} % ajusta o espaçamento dos parágrafos do resumo
\begin{resumo}

 \begin{otherlanguage*}{english}
\textbf{Palavras chaves}
   \vspace{\onelineskip}
     MPO;
     Observatório de Projetos;
     Transparência;
     Ontologias;
     Methontology
   \noindent 
 \end{otherlanguage*}

Esta dissertação apresenta a OntoMPO, uma ontologia de domínio desenvolvida com o objetivo de formalizar o Modelo para Observatórios de Projetos (MPO), originalmente proposto como um modelo conceitual para apoiar a análise e o desenvolvimento de observatórios de projetos. Apesar de sua relevância, o MPO é descrito em linguagem natural, o que limita sua precisão semântica, dificulta a interoperabilidade e restringe sua aplicação em contextos computacionais. Para superar essas limitações, esta pesquisa propõe a construção de uma ontologia que explicita, de forma sistemática, os conceitos, relações, atributos e restrições semânticas do modelo. O desenvolvimento da ontologia segue a metodologia \textit{Methontology}, contemplando as etapas de especificação, conceitualização, formalização, integração, implementação e avaliação. A modelagem ontológica é fundamentada em OntoUML, garantindo rigor semântico, e sua implementação é realizada na linguagem OWL. A avaliação da ontologia envolve validação conceitual e verificação lógica, com base em questões de competência, assegurando sua consistência interna e aderência ao modelo original. Como resultado, a OntoMPO estabelece uma base semântica coerente e formalmente definida, contribuindo para a redução de ambiguidades, padronização conceitual e apoio a futuras aplicações no domínio de observatórios de projetos. Do ponto de vista prático, a OntoMPO pode ser utilizada como base semântica para apoiar o desenho de novos observatórios de projetos, bem como para possibilitar a comparação estruturada entre diferentes observatórios.
\end{resumo}

% resumo em inglês
\begin{resumo}[Abstract]
 \begin{otherlanguage*}{english}
   
\textbf{Keywords}
   \vspace{\onelineskip}
     MPO;
     Project Observatory;
     Transparency;
     Ontologies;
     Methontology;
   \noindent 
 \end{otherlanguage*}

This dissertation presents OntoMPO, a domain ontology developed with the aim of formalizing the Model for Project Observatories (MPO) \cite{vieira2022thesis}, originally proposed as a conceptual model to support the analysis and development of project observatories. Despite its relevance, MPO is described in natural language, which limits its semantic precision, hinders interoperability and restricts its application in computational contexts. To overcome these limitations, this research proposes the construction of an ontology that systematically explains the concepts, relationships, attributes and semantic restrictions of the model. The development of the ontology follows the methodology \textit{Methontology} \cite{fernandez1997methontology}, covering the stages of conception, conceptualization, formalization, integration, implementation and evaluation. The ontological modeling is based on OntoUML \cite{ontouml_readthedocs}, ensuring semantic rigor, and its implementation is carried out in the OWL language. The ontology assessment involves conceptual validation and logical verification, based on competence issues, ensuring its internal consistency and adherence to the original model. As a result, OntoMPO establishes a coherent and formally defined semantic basis, contributing to the reduction of ambiguities, standardization and support for future conceptual applications in the field of project observatories. From a practical point of view, OntoMPO can be used as a semantic basis to support the design of new project observatories, as well as to enable a structured comparison between different observatories.
 
\end{resumo}


